%%%%%%%%%%%%%%%%%%%%%%%%%%%%%%%%%%%%%%%%%
% Arsclassica Article
% LaTeX Template
% Version 1.1 (1/8/17)
%
% This template has been downloaded from:
% http://www.LaTeXTemplates.com
%
% Original author:
% Lorenzo Pantieri (http://www.lorenzopantieri.net) with extensive modifications by:
% Vel (vel@latextemplates.com)
%
% License:
% CC BY-NC-SA 3.0 (http://creativecommons.org/licenses/by-nc-sa/3.0/)
%
%%%%%%%%%%%%%%%%%%%%%%%%%%%%%%%%%%%%%%%%%

%----------------------------------------------------------------------------------------
%	PACKAGES AND OTHER DOCUMENT CONFIGURATIONS
%----------------------------------------------------------------------------------------

\documentclass[
10pt, % Main document font size
a4paper, % Paper type, use 'letterpaper' for US Letter paper
oneside, % One page layout (no page indentation)
%twoside, % Two page layout (page indentation for binding and different headers)
headinclude,footinclude, % Extra spacing for the header and footer
BCOR5mm, % Binding correction
]{scrartcl}

\input{structure.tex} % Include the structure.tex file which specified the document structure and layout

\hyphenation{Fortran hy-phen-ation} % Specify custom hyphenation points in words with dashes where you would like hyphenation to occur, or alternatively, don't put any dashes in a word to stop hyphenation altogether

\usepackage{siunitx}
\usepackage[american,RPvoltages]{circuiTikz}
\usepackage{tikz, tikz-cd}
\usepackage{pgfplots}
\usepackage{float}
\usepackage{graphicx}
\usepackage{xcolor}
\usepackage{listings}

\setlength{\parindent}{0pt}

\definecolor{mGreen}{rgb}{0,0.6,0}
\definecolor{mGray}{rgb}{0.5,0.5,0.5}
\definecolor{mPurple}{rgb}{0.58,0,0.82}
\definecolor{backgroundColour}{rgb}{0.95,0.95,0.92}

\newcommand{\listingsttfamily}{\fontfamily{NotoSansMono-TLF}\small}
\lstdefinestyle{CStyle}{
	language=C,
	backgroundcolor=\color{backgroundColour},   
	commentstyle=\color{mGreen},
	stringstyle=\color{mPurple},
	keywordstyle=\color{magenta},
	keywordstyle = [2]{\color{blue}},
	keywordstyle = [3]{\color{blue}},
	numberstyle=\tiny\color{mGray},
	otherkeywords = {uint8_t,uint16_t,_Bool},
	morekeywords = [2]{uint8_t,uitn16_t},
	morekeywords = [3]{.},
	basicstyle=\listingsttfamily,
	breakatwhitespace=false,         
	breaklines=true,                 
	captionpos=b,                    
	keepspaces=true,                 
	numbers=left,                    
	numbersep=5pt,                  
	showspaces=false,                
	showstringspaces=false,
	showtabs=false,                  
	tabsize=1,
}
%----------------------------------------------------------------------------------------
%	TITLE AND AUTHOR(S)
%----------------------------------------------------------------------------------------

\title{\normalfont\spacedallcaps{Coeficiente de temperatura de las resistencias}} % The article title

%\subtitle{Subtitle} % Uncomment to display a subtitle

\author{Rodrigo Castillejos} % The article author(s) - author affiliations need to be specified in the AUTHOR AFFILIATIONS block

\date{} % An optional date to appear under the author(s)

%----------------------------------------------------------------------------------------

\begin{document}

%----------------------------------------------------------------------------------------
%	HEADERS
%----------------------------------------------------------------------------------------

\renewcommand{\sectionmark}[1]{\markright{\spacedlowsmallcaps{#1}}} % The header for all pages (oneside) or for even pages (twoside)
%\renewcommand{\subsectionmark}[1]{\markright{\thesubsection~#1}} % Uncomment when using the twoside option - this modifies the header on odd pages
\lehead{\mbox{\llap{\small\thepage\kern1em\color{halfgray} \vline}\color{halfgray}\hspace{0.5em}\rightmark\hfil}} % The header style

\pagestyle{scrheadings} % Enable the headers specified in this block

%----------------------------------------------------------------------------------------
%	TABLE OF CONTENTS & LISTS OF FIGURES AND TABLES
%----------------------------------------------------------------------------------------

\maketitle % Print the title/author/date block

\setcounter{tocdepth}{2} % Set the depth of the table of contents to show sections and subsections only

\tableofcontents % Print the table of contents

\listoffigures % Print the list of figures

\listoftables % Print the list of tables

%----------------------------------------------------------------------------------------
%	ABSTRACT
%----------------------------------------------------------------------------------------

\newpage
\section*{Resumen} % This section will not appear in the table of contents due to the star (\section*)

Este documento presenta el diseño e implementación de un sistema de medición de corriente basado en el monitor de potencia, corriente y voltaje ultra preciso de 16 bits con interfaz $I^2C$, el INA236 de Texas Instruments. Se abordará la creación de una librería en lenguaje C para la correcta comunicación entre el INA236 y una placa de desarrollo de STMicroelectronics (Nucleo 64) equipada con el microcontrolador STM32F303XX, aprovechando sus periféricos $I^2C$ y sus capacidades de procesamiento para aplicaciones de adquisición de datos. La librería permitirá no solo la lectura de corriente, voltaje y potencia, sino también la configuración dinámica del valor del resistor shunt , lo que facilita la adaptación del sistema a diferentes rangos de corriente sin modificar el hardware. \\
El diseño incluirá el cálculo del resistor de shunt además de la configuración de los registros internos del INA236 para ajustar las resoluciones y los rangos dinámicos. Finalmente, se desarrollará un diseño de PCB en Altium Designer. Este trabajo integra hardware, software y diseño electrónico para ofrecer una solución robusta en aplicaciones de monitoreo energético.

%----------------------------------------------------------------------------------------
%	AUTHOR AFFILIATIONS
%----------------------------------------------------------------------------------------

%\let\thefootnote\relax\footnotetext{* \textit{Department of Biology, University of Examples, London, United Kingdom}}

%\let\thefootnote\relax\footnotetext{\textsuperscript{1} \textit{Department of Chemistry, University of Examples, London, United Kingdom}}

%----------------------------------------------------------------------------------------

\newpage % Start the article content on the second page, remove this if you have a longer abstract that goes onto the second page

%----------------------------------------------------------------------------------------
%	INTRODUCTION
%----------------------------------------------------------------------------------------

\section{Introducción}
Entender el coeficiente de temperatura de la resistencia (TCR) es crítico para asegurar la estabilidad del resistor en circuitos electrponicos. El coeficiente de temperatura, medido en \textit{partes por millón por grado centígrado (ppm/°C)}, define qué tanto la resistencia de un resistor cambia con la temperatura. Un TCR alto puede generar variaciones significativas en la resistencia, lo que afecta el rendimiento del circuito, mientras que las resistencias con TCR bajo ofrecen mejor estabilidad para aplicaciones de precisión.

\section{¿Qué es el coeficiente de temperatura?}
El coeficiente de temperatura de la reistencia (\textit{Temperature Coefficient of Resistance} (TCR)), a veces referido como coeficiente de temperatura de la resistencia (\textit{resistance temperature coefficient (RTC)}), indica qué tanto la resistencia incrementará o o reducirá por cada $\qty{1}{\degreeCelsius}$ de variación en la temperatura. El efecto de esta variación en la resistencia es reversible en cuanto la temperatura regresa a la temperatura de referencia. \\

Esta variación en la resistencia debido a la temperatura se mide en $ppm/ ^{\circ} C$ que varía ampliamente entre diferentes materiales. Por ejemplo, la aleación de manganeso y cobre tiene un TCR $< 20 ppm/ ^{\circ} C$ (para un rango de temperatura de \qty{20}{\degreeCelsius} a \qty{60}{\degreeCelsius}), mientras que el cobre usado en las terminaciones es aproximadamente \qty{3900}{ppm / \degreeCelsius}

\section{¿Cómo se mide el TCR?}
El rendimiento del TCR según el método 304 de la normal \textbf{MIL-STD-202} se basa en el cambio de resistencia con una temperatura de referencia de \qty{25}{\degreeCelsius}. Se modifica la temperatrura y se permite que el dispositivo bajo prueba alcance el equilibrio antes de medir el valor de la resistencia. La diferencia es usada para determinar el TCR.\\

\textbf{Coeficiente de temperatura (\%)}
\begin{align}
	\frac{R_2 - R_1}{R_1(t_2 - t_1)} \times 100
\end{align}

\textbf{Coeficiente de temperatura (ppm)}
\begin{align}
	\frac{R_2 - R_1}{R_1(t_2 - t_1)} \times 1\,000\,000
\end{align}

Donde:

\begin{description}
	\item[$R_1$ =] Valor de la resistencia a la temperatura de referencia
	\item[$R_2$ = ] Valor de la temperatura a la temperatura de operación
	\item[$t_1$ = ] Temperatura de referencia \qty{25}{\degreeCelsius}
	\item[$t_2$ = ] Temperatura de operación
\end{description}


\begin{align}
	V(t) = V_0 e^{-t/\tau} \nonumber
\end{align}

\begin{align}
	i_R(t) = \frac{v(t)}{R} = \frac{V0}{R} e^{-t/\tau} \nonumber
\end{align}
\begin{align}
	i_C(t) = C\frac{dv(t)}{dt} = -\frac{1}{\tau}V_0 C \cdot e^{-t/\tau} = -\frac{V_0}{R} e^{-t/\tau} \nonumber
\end{align}

\begin{align}
	i_R(t=0) = \frac{V_0}{R} \nonumber \\
	i_C(t=0) = -\frac{V_0}{R} \nonumber
\end{align}

\begin{align}
	\tau = RC = (\qty{330}{\ohm})(\qty{4.7}{\micro \farad}) = \qty{1.551}{\milli\second}
\end{align}
%----------------------------------------------------------------------------------------
%	BIBLIOGRAPHY
%----------------------------------------------------------------------------------------

%\renewcommand{\refname}{\spacedlowsmallcaps{References}} % For modifying the bibliography heading

%\bibliographystyle{unsrt}

%\bibliography{sample.bib} % The file containing the bibliography

%----------------------------------------------------------------------------------------

\end{document}